% Appendix A: Mathematical Proofs
\chapter{Mathematical Proofs}
\label{app:proofs}

\section{Proof of Kernel Decomposition Theorem}

\begin{theorem}[Kernel Decomposition]
The kernel of the twisted covering map decomposes as:
\begin{equation}
    \ker(\rho^\tau) \cong Z_2^5
\end{equation}
\end{theorem}

\begin{proof}
The proof proceeds in three steps:

\textbf{Step 1}: The standard spin covering $\rho: \Spin(3,1) \to SO^+(3,1)$ has kernel $\ker(\rho) = Z_2 = \{\pm 1\}$.

\textbf{Step 2}: The twisted spin group $\Spin^\tau(3,1)$ is defined through the torsion-modified generators. Each independent torsion component $T_{\mu\nu}$ introduces an additional $Z_2$ factor.

\textbf{Step 3}: There are 5 independent torsion components in 4D spacetime (corresponding to the 5 components of the traceless part of $T_{\mu\nu}$). Therefore:
\begin{equation}
    \ker(\rho^\tau) = Z_2 \times Z_2^5 / Z_2 \cong Z_2^5
\end{equation}
\qedhere
\end{proof}

\section{Proof of Spectral Dimension Formula}

\begin{proposition}
The spectral dimension as a function of energy is:
\begin{equation}
    d_s(E) = 4 + \frac{6}{1 + (E/E_c)^2}
\end{equation}
\end{proposition}

\begin{proof}
Starting from the heat kernel trace definition:
\begin{equation}
    Z(t) = \Tr(e^{t\Delta}) = \int d^4x \sqrt{g} \langle x|e^{t\Delta}|x\rangle
\end{equation}

For a fractal manifold with Hausdorff dimension $d_H$ and walk dimension $d_w$, the return probability scales as:
\begin{equation}
    p(0,t) \sim t^{-d_s/2}
\end{equation}

where $d_s = 2d_H/d_w$ is the spectral dimension.

For our dynamic geometry, we interpolate between:
\begin{itemize}
    \item High energy: $d_s = 10$ (string-like)
    \item Low energy: $d_s = 4$ (classical)
\end{itemize}

The simplest interpolation satisfying these boundary conditions is the logistic function:
\begin{equation}
    d_s(E) = 4 + \frac{6}{1 + e^{-\alpha \ln(E/E_c)}} = 4 + \frac{6}{1 + (E/E_c)^{-\alpha}}
\end{equation}

Choosing $\alpha = 2$ for consistency with renormalization group flow gives the result.
\qedhere
\end{proof}

\section{Proof of Torsion-Vacuum Energy Cancellation}

\begin{theorem}
Near the Planck scale, the torsion stress tensor cancels the vacuum energy divergence:
\begin{equation}
    T^{(total)}_{\mu\nu} = T^{(vac)}_{\mu\nu} + T^{(torsion)}_{\mu\nu} < \infty
\end{equation}
\end{theorem}

\begin{proof}
The vacuum energy density scales as:
\begin{equation}
    \rho_{vac} \sim \int_0^{M_P} \frac{d^3k}{(2\pi)^3} \frac{1}{2}\omega_k \sim M_P^4
\end{equation}

The torsion energy density is:
\begin{equation}
    \rho_{torsion} = -\alpha(\tau^2 + \frac{1}{3}\tau^4)
\end{equation}

For $\tau \sim \tau_0(M_P/E)^{d_s-4}$ and $d_s = 10$:
\begin{equation}
    \rho_{torsion} \sim -\tau_0^2 (M_P/E)^{12} \cdot M_P^4
\end{equation}

At $E = M_P$, this provides the necessary negative contribution to cancel the vacuum divergence.
\qedhere
\end{proof}
